\documentclass[11pt,letterpaper]{article}
\usepackage[T1]{fontenc}
\usepackage{lmodern,amsmath,amssymb,amsthm,microtype}
\usepackage[margin=1in]{geometry}
\usepackage[hidelinks]{hyperref}
\providecommand{\tightlist}{\setlength{\itemsep}{0pt}\setlength{\parskip}{0pt}}
\DeclareUnicodeCharacter{2264}{\ensuremath{\le}}
\DeclareUnicodeCharacter{2265}{\ensuremath{\ge}}
\DeclareUnicodeCharacter{2260}{\ensuremath{\ne}}
\DeclareUnicodeCharacter{2212}{\ensuremath{-}}

\DeclareUnicodeCharacter{2211}{\ensuremath{\sum}}

\title{Only the final eighth-weight floor survives below 150 for nine separated subset sums}
\author{Aafi Khan Malek\\ \small Independent researcher, Ahmedabad, India\\
\small\href{mailto:aafimalek2032@gmail.com}{aafimalek2032@gmail.com}}
\date{28 November 2026}
\begin{document}
\maketitle
\textbf{Theorem.} Let nine sorted positive real weights
have pairwise unit-separated subset sums. If \(149\le x_9<150\), then
\[
 148\le x_8<149,\qquad
 1\le x_9-x_8\le 1+\frac{x_9-148}{39}<\frac{41}{39}.
\]
We use the existing
\href{https://github.com/Aafimalek/novel-discoveries/tree/aee69772bd5c97a7f956f92e2ad5d1cacdac0a2a/nine-separated-subset-sums/verification/rank-cuts-below-150/penultimate-cap-135}{exact exclusion of eighth weights below 136}
and the previous Harper-cut argument. Mixed rounding and rank-swap
counting then eliminate twelve additional floors without chamber
enumeration. A narrow remaining case is still open.

\section{Mixed rounding forbids a saturated coordinate-eight
edge}\label{mixed-rounding-forbids-a-saturated-coordinate-eight-edge}

Suppose \(t>128\) and an edge changing bit eight has rank length \(t\).
After complementation, let its lower rank be \(\ell\). It is impossible
if \(257-t\le\ell\le177\). Use integer weights
\(a_i=\lfloor x_i\rfloor\) for \(i\le8\) and \(a_9=t+D\). Thus \(a_9\)
is 148 or 149, and \(a_9-a_8=\lfloor x_9-x_8\rfloor\). Write \(Q\) for
the total real weight of the lower endpoint of the edge.

The \(t+1\) ranked scores in \([Q,Q+x_8]\) have total spacing excess
\(x_8-t<1\). Any positive in-interval difference has floor equal to its
rank difference and fractional part at most \(x_8-t\). Deleting bit
eight identifies exactly the two endpoints: any other pair of lifts of
one projection would be \(x_8\) apart and would have to use the same
endpoints. Hence the projection contains \(t>128\) cube vertices. For
every other bit, its perfect matching has an edge entirely in this
projection.

For a bit \(i\le7\), a projected edge can lift to an in-interval
difference \(x_i\) or \(x_8-x_i\); the sum \(x_8+x_i\) cannot fit. Their
fractional parts force \(\{x_i\}\le\{x_8\}\): if the inequality failed
both differences would have excessive fractional part. Consequently in
either type of lift the floor of the positive real difference equals the
difference of the integer scores from the \(a_i\). A projected edge
changing bit nine instead must lift with opposite eighth bits: the only
positive difference that can fit is \(x_9-x_8\). It too has floor equal
to the integer score difference, by the mixed choice of \(a_9\).

The projected vertices are the scores in the slab \([Q-x_8,Q+x_8]\). In
each fixed-bit-nine sheet this slab is connected using differing bits
from the first seven coordinates. While both an addition and a deletion
remain, add when the current score is at most \(Q\), otherwise delete;
all steps have size at most \(x_8\). When only one direction remains,
the path is monotone to the target. The projected bit-nine matching edge
joins the sheets. Lifting paths is possible because the only vertex with
two lifts is the projection of the edge endpoints, whose rank difference
is also exactly their integer score difference. Thus rank minus integer
score is constant on the entire edge interval.

The reflected interval gives the same local constancy and meets this
interval because \(\ell\ge257-t\). The connected union is centered at
half the total real weight and has half-width at least
\((513-2\ell)/2\ge79.5\) by rank separation. Any equal pair of integer
subset sums can first have their common coordinates removed, and then be
translated by the same subset of the unused coordinates. Choose a prefix
closest to half the unused integer total; consecutive prefix scores
differ by at most \(a_9\). The translated pair lies within \(a_9/2\) of
the integer center. Real rounding errors are all nonnegative. If
\(a_9=149\) their total is less than 9, and if \(a_9=148\) their total
is less than 10. In either case both real subset sums therefore lie
strictly within 79 of the real center. The constant rank-minus-integer
score on the union excludes such a collision. Integer subset-sum
injectivity would contradict
\href{https://www.cecm.sfu.ca/~pborwein/PAPERS/P170.pdf}{Borwein and
Mossinghoff's integer optimum of 161 for nine weights}, Table 5. This
proves the forbidden band, without any assumption that \(D\) equals
\(149-t\).

\section{Extend every low-side rank run by
one}\label{extend-every-low-side-rank-run-by-one}

Put \(M=x_9\), \(d=M-x_8\), \(t=\lfloor x_8\rfloor\), and
\(D=\lfloor d\rfloor\le149-t\). Unit separation gives
\(x_8\le M-1<149\), so \(136\le t\le148\) after the previous bound.

At each of the eight low \href{../README.md}{Harper cuts}
\(k=158,159,161,162,168,169,171,172\), the original length-149 forbidden
band and the boundary size 148 give all boundary ranks \(k+1\) through
\(k+148\). Consider the upper ranks \(j=k+t,\ldots,k+148\). Coordinates
below eight have weight less than \(t\), so they cannot give a rank jump
of at least \(t\). To reach \(j=k+t\) via bit eight would require an
exact length-\(t\) edge starting at \(k\), which the mixed-rounding band
excludes. For larger \(j\) bit eight is already too short. Hence each of
these boundary ranks is reached via bit nine. A length-149 bit-nine edge
would start at \(j-149\) in the original forbidden band, so each lower
endpoint is at least \(j-148\). The lower endpoints strictly increase
with \(j\) because their scores are obtained by subtracting the same
\(M\), and the last one is exactly \(k\). Every lower rank is therefore
exactly \(j-148\).

Thus the prescribed starts extend to every rank in \([k+t-148,k]\) at
each cut. This is one start earlier than the
\href{../hall-rank-runs/}{previous rank-run lemma}, since the saturated
bit-eight alternative is forbidden even without a floor-gap equality.
Write \(S\) for the union of these low starts.

\section{Swap capacity and the twelve eliminated
floors}\label{swap-capacity-and-the-twelve-eliminated-floors}

For a forced bit-nine edge from \(s\) to \(s+148\), if bit eight is
absent at the lower endpoint, swap nine for eight at the upper endpoint.
The rank decreases by \(h\in[1,D]\) and \(s-h\notin S\), since the
swapped upper vertex omits nine. If bit eight is present at the lower
endpoint, swap eight for nine there instead. The rank increases by
\(h\in[1,D]\) and \(s+h\notin S\), since the swapped lower vertex now
contains nine. Distinct starts of the same type have distinct partners
because swapping the bits is an involution.

For any consecutive rank span of length \(L\) with \(H\) holes in \(S\),
its \(L-H\) starts have at most \(D+H\) left partner slots and \(D+H\)
right partner slots, even if one ignores the other forced starts outside
the span. Therefore \(L\le2D+3H\). Apply this as follows:

\begin{itemize}
\tightlist
\item
  For \(136\le t\le143\), the runs merge into \([t+10,172]\), a full
  span of length \(163-t>2(149-t)\).
\item
  For \(t=144\), the low span \([154,172]\) has length 19 and one hole.
  The available capacity is at most \(2\cdot5+3=13\).
\item
  For \(t=145\), the low span \([155,172]\) has length 18 and two holes.
  The capacity is at most \(2\cdot4+6=14\).
\item
  At \(t=146\), one contiguous run is \([156,162]\), length 7, but
  \(D\le3\).
\item
  At \(t=147\), one contiguous run is \([157,162]\), length 6, but
  \(D\le2\).
\end{itemize}

These contradict the capacity and leave only \(t=148\).

\section{A narrow gap at floor 148}\label{a-narrow-gap-at-floor-148}

Now \(1\le d<2\), so each swap changes rank by exactly one. There are
four consecutive pairs \(k,k+1\) of forced starts, with
\(k=158,161,168,171\). In each pair, start \(k\) must use the left move,
because a right move would land on start \(k+1\). Likewise \(k+1\) must
use the right move. Thus the gap from rank \(k+147\) to rank \(k+148\)
and the gap from rank \(k+1\) to rank \(k+2\) both equal \(d\).
Complementing the first gap shows that the gap from rank \(365-k\) to
rank \(366-k\) also equals \(d\).

A stronger argument uses the coordinate-eight companion of the forced
left move at rank 158. Its rank interval is \([158,305]\), with 148
scores of span \(x_8\). The first point has neither of the top bits; the
last has only bit eight, with its swap mate at rank 306, outside this
interval.

Write \(\varepsilon=d-1\ge0\). The first seven weights have
\((1+d)\)-separated subset sums. Indeed a positive difference \(g\) of
two first-seven scores is at least 1, and the two translates by \(x_8\)
and \(M\) give \(|g-d|\ge1\). Since \(d<2\), the branch \(g\le d-1\) is
impossible, so \(g\ge1+d\).

Call a rank middle if exactly one of the two top bits is present and
outer otherwise. All middle points occur in adjacent swap pairs of score
gap \(d\): a pair is adjacent because \(d<2\). Inside \([158,305]\)
there are \(P\) complete middle pairs and just one partial middle pair,
at the last rank. Each complete pair uses \(\varepsilon\) of the
interval's excess over unit spacing.

For a maximal run of \(L\) outer ranks, take \(\lfloor(L-1)/2\rfloor\)
triples of consecutive ranks, sharing one endpoint between successive
triples so their gaps are disjoint. Among the three ranks in any triple
two have the same top-bit pattern, either both empty or both full. Their
score difference is at least \(1+d=2+\varepsilon\) by the first-seven
separation. Thus the triple uses at least \(\varepsilon\) excess in its
two gaps. Let \(F\) count all these triples, and \(r\) count the outer
runs. Then \(r\le P+1\).

At least five runs are singleton: ranks 158, 161, 171, 196, and 206.
This follows from the three-rank patterns forced by the four consecutive
pairs. For each first start \(k\), ranks \(k,k+1,k+2\) have patterns
\(00,10,01\); at \(k+147,k+148,k+149\) the patterns are \(10,01,11\).
Complementing the latter gives \(00,10,01\) at \(364-k,365-k,366-k\).
These patterns place a middle rank on either side of each listed
singleton, unless the side is the interval boundary. The exact replay
checks the five sites.

Ordinarily a run contributes at most \(2\lfloor(L-1)/2\rfloor+2\) ranks;
each known singleton improves that by one. Consequently \[
 148\le 2P+1+2F+2r-5\le4(P+F)-2,
\] so \(P+F\ge38\). Complete middle-pair gaps and outer-triple gaps are
disjoint. The interval has excess \(x_8-147=M-148-\varepsilon\), whence
\[
 38\varepsilon\le M-148-\varepsilon,
 \qquad d=1+\varepsilon\le1+(M-148)/39<41/39.
\]

\texttt{python3\ rank\_swap\_replay.py} independently checks the exact
interval arithmetic and also the necessary bipartite matching conditions
on both low and reflected starts. The orbit and cap replays in the
preceding work are separate from this handwritten proof. This is not a
Lean theorem and does not exclude floor 148.

\end{document}
